% Research Aim and Hypothesis
\section{Research Aim and Hypothesis}

\subsection{Research Aim}

The overarching aim of this research is to develop, implement, and validate a Hybrid Sparse-Sampling
Architecture for Reasoning Video Segmentation. This architecture seeks to decouple the heavy semantic
processing from the spatial tracking task.

Specifically, the aim is to create a system that utilizes Deep Learning
(\cite[VideoLISA]{bai_one_2024}/\cite[SAM 2]{ravi_sam_2024})
only on sparse keyframes to establish semantic understanding, while employing lightweight Classical
Interpolation (Kriging) and Trajectory Forecasting (\cite[LSTM]{hochreiter_long_1997}) to reconstruct
the intermediate frames.
By doing so, the research aims to prove that it is possible to achieve State-of-the-Art (SOTA)
accuracy and stability on resource-constrained hardware, effectively solving the "Stability-Efficiency Trade-off."

\subsection{Research Hypothesis}

The study is guided by the following core hypothesis and sub-hypotheses:

\textbf{Core Hypothesis:}

\textit{It is hypothesized that a hybrid segmentation framework, which combines sparse semantic
    keyframe extraction with geostatistical interpolation and LSTM-based trajectory forecasting, will achieve
    higher temporal stability and computational efficiency than a purely deep-learning-based dense sampling approach,
    without a statistically significant reduction in semantic accuracy.}

\textbf{Sub-Hypothesis 1 (Efficiency):}

\textit{By reducing the frequency of deep learning inference to 20\% of the video frames (Sparse Sampling)
    and utilizing interpolation for the remaining 80\%, the system will achieve a processing speed increase of at
    least 200\% (3x speedup) compared to the baseline \cite[VideoLISA]{bai_one_2024} model.}

\textbf{Sub-Hypothesis 2 (Stability):}

\textit{The application of Geostatistical Interpolation (Universal Kriging) on mask contours will result in
    a lower "Flow-Warping Error" (a metric of temporal jitter) compared to masks generated independently by neural
    networks, due to the variance-minimizing properties of the Kriging estimator.}

\textbf{Sub-Hypothesis 3 (Occlusion):}

\textit{An LSTM-based motion predictor, trained on object centroid trajectories, will maintain higher
    Intersection-over-Union (IoU) scores during periods of full occlusion (5+ frames) compared to the standard
    visual attention mechanisms of Transformers, as the LSTM explicitly models momentum and velocity.}